% ============================================================================
\section{Algoritmos que sostienen el catálogo}
\label{sec:catalogo}
% ============================================================================

La resolución y la optimización \emph{suponen} propiedades del catálogo:
que el empaque de cada producto sea correcto, que el nivel de marca esté
ordenado por precio dentro de cada concepto, que haya exactamente un producto
representante por celda y que los vectores correspondan al texto vigente.
Ninguna de esas propiedades se cumple sola: las garantizan los algoritmos de
esta sección, que se ejecutan al curar el catálogo y no al atender una
solicitud.

% ---------------------------------------------------------------------------
\subsection{Extracción del empaque desde el nombre del producto}
\label{sec:cat-empaque}
% ---------------------------------------------------------------------------

Todo el arreglo de la cantidad (\cref{sec:nor-cantidad}) depende de
$\emp(p)$. Con $\emp\equiv1$ ---el estado inicial del catálogo, en el que los
260 productos tenían empaque 1---, la fórmula $\techo{q/\emp}$ no hace nada.
El empaque está escrito en el nombre del producto, y la dificultad no es
encontrar números sino decidir \emph{cuál cuenta}: «Témpera gouache de 12~ml
x 12 unidades» tiene dos doces y solo el segundo multiplica; «Cinta 18mm X5
Metros X10 Unds» tiene tres números y solo el último cuenta.

\begin{definicion}[Regla del empaque]
\label{def:regla-empaque}
Sea $\bar n$ el nombre en forma plana y sea $D(c)$ el conjunto de
discriminantes del concepto. Se recorren las coincidencias de
\[
  (\text{inicio}\mid\text{separador})\ (\texttt{x}\mid\texttt{por})\ N\ (\texttt{.}d^+)?\ u?
\]
con $N$ entero y $u$ una palabra opcional, y cada coincidencia se descarta si:
\begin{enumerate*}[label=(\roman*)]
  \item tiene parte decimal («30,5 x 30,5 cm»: nadie empaca media unidad);
  \item $u$ es una unidad de \emph{medida} (gramos, ml, cm, metros, \dots); o
  \item $u\in\{\text{hojas}\}$ y $\text{hojas}\in D(c)$.
\end{enumerate*}
Una palabra $u$ que no cuenta ni mide se ignora («x100 Gigante» sigue siendo
100). Entre las coincidencias que sobreviven gana la \emph{mayor}, y compite
también con la forma escrita «paquete/caja/set/juego/bolsa/blíster/kit de
$N$». Si no queda ninguna y las hojas no son atributo, se acepta un conteo
suelto «$N$ hojas/pliegos/unidades». En otro caso, $\emp(p)=1$.
\end{definicion}

La regla, en una línea, es: \emph{«x$N$» cuenta empaques salvo que lo siga
una unidad de medida}. La excepción (iii) resuelve el único caso que la
sintaxis sola no puede resolver. Cien hojas de un cuaderno son una
\emph{propiedad} del cuaderno; cinco hojas de acetato son \emph{cinco cosas}.
La misma palabra significa cosas opuestas según el producto, y por eso lo
decide el concepto: si declara \cod{hojas} como discriminante, las hojas
describen y no cuentan. Elegir la mayor entre varias coincidencias recoge el
caso de un «x2» de presentación junto a un «x24» de caja, en el que el
empaque real es el grande. La complejidad es lineal en la longitud del nombre.

% ---------------------------------------------------------------------------
\subsection{Nivelación: etiquetas de nivel monótonas en precio}
\label{sec:cat-nivelar}
% ---------------------------------------------------------------------------

El nivel de marca es uno de los objetivos del optimizador
(\cref{sec:opt-objetivos}). Si dentro de un concepto el producto etiquetado
premium resulta más barato que el económico, el acudiente que elige
«económico» paga más que el que elige «premium»: la promesa de los kits queda
invertida, y nada en el resto del sistema lo detecta.

\begin{definicion}[Escalera y nivelación]
Una \emph{escalera} es el conjunto $E=\{p_1,\dots,p_k\}$ de productos
\emph{representantes} activos de un mismo par (concepto, personaje), cada uno
con un precio $\precio(p_i)$ y una etiqueta $\lambda(p_i)\in\{0,1,2\}$. Una
\emph{nivelación} de $E$ es una etiqueta nueva $\lambda'$ tal que:
\begin{enumerate}[label=(N\arabic*)]
  \item\label{niv:multi} el multiconjunto de etiquetas se conserva:
        $\{\!\{\lambda'(p)\}\!\}=\{\!\{\lambda(p)\}\!\}$;
  \item\label{niv:mono} $\lambda'$ es monótona en el orden total
        $p\sqsubseteq p'\iff(\precio(p),\mathrm{id}(p))\leq_{\mathrm{lex}}(\precio(p'),\mathrm{id}(p'))$.
\end{enumerate}
\end{definicion}

\begin{proposicion}[Existencia y unicidad]
\label{prop:nivelar}
Toda escalera tiene exactamente una nivelación, que se obtiene ordenando los
productos por $\sqsubseteq$, ordenando las etiquetas de menor a mayor y
asignándolas en ese orden. Se calcula en $O(k\log k)$.
\end{proposicion}

\begin{proof}
\emph{Existencia}: la asignación descrita conserva el multiconjunto, porque
reparte exactamente las mismas etiquetas, y es monótona por construcción.
\emph{Unicidad}: sea $\lambda'$ una nivelación y sean $p_{(1)}\sqsubset\cdots
\sqsubset p_{(k)}$ los productos ordenados. Por \ref{niv:mono},
$\lambda'(p_{(1)})\leq\cdots\leq\lambda'(p_{(k)})$, así que la secuencia
$(\lambda'(p_{(i)}))_i$ es la única ordenación no decreciente de su
multiconjunto, que por \ref{niv:multi} es el de $\lambda$. Dos ordenaciones no
decrecientes de un mismo multiconjunto son iguales. El costo lo dominan los
dos ordenamientos.
\end{proof}

La reparación es deliberadamente mínima. No inventa niveles ni reparte por
terciles: cuántos productos económicos, confiables o premium tiene cada
concepto no cambia, solo cambia qué producto lleva cada etiqueta. Además,
solo se modifican las escaleras cuyo orden estaba roto, de modo que la
diferencia muestra exactamente dónde estaba el defecto. Dos detalles completan
el algoritmo:
\begin{itemize}
  \item \textbf{Escaleras paralelas por personaje.} El cuaderno con personaje
        económico compite con el cuaderno con personaje premium, no con el
        cuaderno sin personaje. Por eso la escalera es por par (concepto,
        personaje).
  \item \textbf{Personaje solitario.} Un producto que es el único de su
        escalera de personaje no tiene con qué compararse. Su etiqueta se lee
        entonces de la \emph{banda de precios} de los productos de marca del
        mismo concepto: es la etiqueta del producto de marca más caro cuyo
        precio no supera el suyo, o la del más barato si todos lo superan. El
        caso que motivó la regla fue un planeador con personaje a \$34.900
        marcado como económico, cuando la agenda de marca premium valía
        \$26.600.
\end{itemize}

% ---------------------------------------------------------------------------
\subsection{Un representante por celda}
\label{sec:cat-representantes}
% ---------------------------------------------------------------------------

La revisión del catálogo puede aprobar varias referencias para la misma celda
(concepto, nivel, personaje). La resolución, sin embargo, ofrece una opción por
nivel (\cref{sec:res-niveles}), y la nivelación opera solo sobre
representantes. El algoritmo de representantes garantiza el invariante
\[
  \forall\,(c,\ell,\mathit{per}):\quad
  \card{\{p\in\Prod : c(p)=c,\ \nivel(p)=\ell,\ \mathrm{per}(p)=\mathit{per},\
  p \text{ es representante}\}}\leq 1,
\]
con igualdad cuando la celda tiene algún producto. Para elegir el
representante se prefiere el producto semilla de la celda, si existe, y si no
el más barato. Los demás quedan como profundidad de compra detrás del
representante. El orden de ejecución al exportar la curaduría es: primero
representantes, después la nivelación (que opera sobre ellos) y por último una
recarga que comprueba que lo exportado se puede leer.

% ---------------------------------------------------------------------------
\subsection{El catálogo como caché progresiva}
\label{sec:cat-cache}
% ---------------------------------------------------------------------------

El catálogo es propio y curado, y solo contiene lo que el sistema puede
conseguir. Una tienda externa interviene en un único momento: cuando llega la
lista de un colegio que nadie ha revisado y un renglón no encuentra nada en
casa. La consecuencia es que el costo baja solo: la primera lista de un
colegio nuevo consulta tiendas por sus renglones raros; al aprobarse, esos
productos quedan en el catálogo, y la siguiente lista ---de ese colegio o de
otro--- los encuentra en casa. Una lista revisada no vuelve a consultar
tiendas, y la estructura del código lo garantiza: el camino que sale a una
tienda solo es alcanzable desde una versión de lista sin revisar.

\paragraph{Búsqueda de candidatos en tiendas.} Dada la consulta $Q(x)$
(\cref{def:cabeza}) y su ancla (la cabeza), el procedimiento es:
\begin{enumerate}
  \item Consultar cada tienda y conservar solo las ofertas disponibles, con
        imagen y con precio en $[\$1.000,\ \$400.000]$. Ese filtro grueso
        descarta las garrafas industriales y los lotes institucionales que
        las tiendas mezclan con el artículo suelto.
  \item Exigir que el nombre de la oferta contenga los primeros cinco
        caracteres del ancla. Es un filtro burdo, pero resuelve el fallo más
        común: pedir «aserrín colores surtidos» devuelve velas y un trapero,
        porque el buscador de la tienda se queda con «colores surtidos» y
        suelta el sustantivo.
  \item Ordenar por precio y, si quedan más de tres, tomar las ofertas en las
        posiciones $\piso{N/6}$, $\piso{N/2}$ y $\piso{5N/6}$: una barata, una
        de en medio y una cara. Se evitan los extremos porque la más barata
        suele ser un accesorio suelto y la más cara un empaque institucional.
\end{enumerate}
El resultado son \emph{propuestas}, nunca productos. Una persona las aprueba o
las rechaza, porque lo que falla en una búsqueda por texto es el significado,
y contra eso ningún umbral sirve: pedirle «punzón punta metálica» a una tienda
devuelve lápices correctores de punta fina, que son productos reales, con
foto, existencias y un precio razonable. Una persona lo ve en dos segundos; un
puntaje de similitud no lo ve nunca. Se proponen tres y no diez porque quien
revisa mira las dos primeras y decide; una cola de revisión larga se abandona,
y una cola abandonada es peor que no tenerla, porque aparenta un control que
nadie ejerce.

\paragraph{Validación de identificadores.} Un código de barras de tienda solo
se guarda si tiene la forma de un EAN real (8, 12 o 13 dígitos). Los de 12
dígitos se normalizan a 13. Un campo con un SKU interno o vacío se descarta,
para no ensuciar el catálogo con una referencia que no sirve para casar
productos.

\begin{evolucion}[: catálogo]
\begin{itemize}
  \item \textbf{Empaque.} De $\emp\equiv1$ en todos los productos a la regla
        de la \cref{def:regla-empaque}. La excepción de las hojas, que depende
        del concepto, se agregó cuando se vio que la misma palabra contaba
        unidades en unos conceptos y describía en otros.
  \item \textbf{Nivelación.} Ninguna prueba comparaba precio contra nivel, y
        varias escaleras estaban invertidas. La nivelación por permutación y
        la prueba que la verifica cerraron el defecto. Después se agregaron
        las escaleras por personaje y la banda de marca para los productos
        solitarios.
  \item \textbf{Representantes.} Al permitir varias referencias por celda,
        permutar etiquetas contra todo el relleno podía quitarle el nivel al
        representante y dárselo a una alternativa que el kit ni siquiera
        ofrecía. La nivelación pasó a operar solo sobre representantes.
  \item \textbf{Huella de los vectores.} Al cambiar el formato del texto
        embebido, los vectores quedaron marcados como desactualizados aunque
        su contenido no hubiera cambiado, y se recalcularon 666. La huella
        hace explícito \emph{de qué} texto salió cada vector.
\end{itemize}
\end{evolucion}
